\documentclass{fiche}
\metadonnees{15}{Le pays des aveugles}{Bloc 3 — La ville, le travail, la science}

\begin{document}
\entete

\objectifs{%
\textbf{Langue} : comparatifs et superlatifs ; \emph{as\ldots as} ; \emph{the more\ldots the
more} ; thème de bilan.\par
\textbf{Texte} : H.\,G. Wells, \og The Country of the Blind \fg (1904).\par
\textbf{Durée} : 1\,h\,30 — exercices 35 min, texte et questions 25 min, version 30 min.}

%% ===================================================================
\exercice[12 min]{Comparatifs et superlatifs}

\rappel{\textbf{Rappel.} Adjectif d'une syllabe (et de deux en \emph{-y}) : \emph{-er},
\emph{-est}. Adjectif plus long : \emph{more}, \emph{most}. Irréguliers : \emph{good},
\emph{bad}, \emph{far}.}

\consigne{Mettez l'adjectif entre parenthèses à la forme qui convient.}

\begin{anglais}
\begin{enumerate}[label=\arabic*.,itemsep=2.2mm,leftmargin=*]
  \item The place seemed \emph{(large)} \trou{larger} as he drew near to it.
  \item The plasterings looked \emph{(queer)} \trou{queerer} than he had expected.
  \item His voice seemed \emph{(coarse)} \trou{coarser} than theirs.
  \item Their notes were \emph{(soft)} \trou{softer} than his.
  \item Perhaps he will grow \emph{(fine)} \trou{finer}.
  \item It was the \emph{(strange)} \trou{strangest} encounter of his life.
  \item Bogota is far \emph{(big)} \trou{bigger} than this village.
  \item This is the \emph{(remote)} \trou{remotest \emph{ou} most remote} valley in the
    Andes.
  \item He was \emph{(good)} \trou{better} at climbing than at arguing.
  \item The valley is \emph{(far)} \trou{farther \emph{ou} further} from Bogota than he
    thought.
  \item That was the \emph{(bad)} \trou{worst} day of the whole journey.
  \item Their hands were \emph{(sensitive)} \trou{more sensitive} than his eyes.
  \item She had the \emph{(sweet)} \trou{sweetest} face in the village.
  \item He found it \emph{(difficult)} \trou{more difficult} than he had anticipated.
\end{enumerate}
\end{anglais}

\note[Trois pièges d'orthographe]{Consonne finale doublée après voyelle brève accentuée :
\emph{big $\rightarrow$ bigger}, \emph{hot $\rightarrow$ hottest}. \emph{-y} devient
\emph{-i} : \emph{happy $\rightarrow$ happier}. Et le superlatif prend toujours \emph{the},
même avec \emph{most} : \emph{the most sensitive}.}

\note[\emph{Farther} ou \emph{further}]{\emph{Farther} pour la distance physique,
\emph{further} pour le reste (\emph{further information}). L'usage les confond de plus en
plus, mais \emph{further} seul peut dire \og supplémentaire \fg.}

%% ===================================================================
\exercice[10 min]{Égalité, proportion, progression}

\rappel{\textbf{Rappel.} \emph{As\ldots as} dit l'égalité, \emph{not so\ldots as} l'inégalité.
\emph{The more\ldots the more} dit la proportion ; \emph{-er and -er} la progression.}

\consigne{Complétez.}

\begin{anglais}
\begin{enumerate}[label=\arabic*.,itemsep=2.2mm,leftmargin=*]
  \item He was \trou{as} rough \trou{as} the rocks that begot him.
  \item Their eyes were \trou{not so} useful \trou{as} their hands.
  \item \trou{The more} he explained, \trou{the less} they understood.
  \item The room grew \trou{darker and darker} as the crowd closed in.
  \item \trou{The nearer} he came, \trou{the stranger} the village looked.
  \item This valley is twice \trou{as} large \trou{as} the other.
  \item He talked \trou{as} confidently \trou{as} a man who enters upon life.
  \item The path became \trou{steeper and steeper}.
  \item \trou{The longer} he stayed, \trou{the more} he doubted his own eyes.
  \item Nothing was \trou{as} useless \trou{as} sight in that country.
\end{enumerate}
\end{anglais}

\note[Construction]{Dans \emph{the more\ldots the more}, les deux membres commencent par
\emph{the} et gardent l'ordre sujet-verbe : \emph{the more he explained, the less they
understood}. Le français dit \og plus\ldots moins \fg{} sans article. Et \emph{twice as
large as} : le multiplicateur se place avant \emph{as}, jamais après.}

%% ===================================================================
\exercice[13 min]{Thème --- bilan du bloc}
\consigne{Traduisez en anglais. Déterminants, relatives, passif et comparaison y sont tous.}

\begin{quote}\itshape
Coketown était une ville de brique rouge, ou de brique qui aurait été rouge si la fumée
l'avait permis. On y comptait plus de cheminées que d'arbres, et la rivière que traversait le
canal était plus noire encore que le ciel. Les hommes et les femmes qui y travaillaient
sortaient tous aux mêmes heures. On dit que la ville la plus riche du comté était aussi la
plus laide. Rien n'y était jamais expliqué : les faits seuls comptaient, et ce qu'on ne
pouvait pas compter n'existait pas.
\end{quote}

\lignes{10}

\corr{\textbf{Proposition.} Coketown was a town of red brick, or of brick which would have
been red if the smoke had allowed it. There were more chimneys than trees in it, and the
river the canal crossed was even blacker than the sky. The men and women who worked there all
went out at the same hours. The richest town in the county is said to have been the ugliest
as well. Nothing was ever explained there: facts alone counted, and what could not be counted
did not exist.}

\note[Quatre pièges]{\og Les hommes et les femmes \fg{} pris en général : pas d'article
(\emph{men and women}), mais ici la relative les détermine, donc \emph{the men and women who
worked there}. \og La rivière que traversait le canal \fg{} : relative complément, le pronom
peut tomber (\emph{the river the canal crossed}). \og On dit que\ldots était \fg{} : passif
d'opinion avec décalage, \emph{is said to have been} (fiche 14). \og Ce qu'on ne pouvait pas
compter \fg{} : \emph{what}, sans antécédent, suffit --- pas de \emph{that which}.}

\note[Comparatifs]{\og plus de cheminées que d'arbres \fg{} : \emph{more\ldots than}, sans
préposition devant le second terme. \og plus noire encore \fg{} : \emph{even blacker}, ou
\emph{blacker still}, mais jamais \emph{more black}.}

%% ===================================================================
\newpage
\section{Texte}

\noindent\small\textit{Nunez, guide de montagne, est tombé dans une vallée fermée des Andes
dont les habitants sont aveugles depuis quinze générations. Il vient d'apercevoir les
premiers d'entre eux.}\normalsize

\begin{texte}
He was sure that they were blind. He was sure that this was the Country of the Blind of which
the legends told. Conviction had sprung upon him, and a sense of great and rather
\gl[enviable]{enviable}{enviable} adventure. The three stood side by side, not looking at him,
but with their ears directed towards him, judging him by his unfamiliar steps. They stood
close together like men a little afraid, and he could see their eyelids closed and
\gl[sunken]{sunken}{enfoncé, creusé} as though the very balls beneath had shrunk away. There
was an expression near \gl[awe]{awe}{l'effroi mêlé de respect} on their faces.

``A man,'' one said, in hardly recognisable Spanish --- ``a man it is --- a man or a spirit
--- coming down from the rocks.''

But Nunez advanced with the confident steps of a youth who enters upon life. All the old
stories of the lost valley and the Country of the Blind had come back to his mind, and through
his thoughts ran this old proverb, as if it were a \gl[a refrain]{refrain}{un refrain} ---

``In the Country of the Blind the One-eyed Man is King.''

And very civilly he gave them greeting. He talked to them and used his eyes.

``Where does he come from, brother Pedro?'' asked one.

``Down out of the rocks.''

``Over the mountains I come,'' said Nunez, ``out of the country beyond there --- where men can
see. From near Bogota, where there are a hundred thousands of people, and where the city
passes out of sight.''

``Sight?'' muttered Pedro. ``Sight?''

``He comes,'' said the second blind man, ``out of the rocks.''

The cloth of their coats Nunez saw was curiously fashioned, each with a different sort of
\gl[stitching]{stitching}{la couture, le point}.

They startled him by a simultaneous movement towards him, each with a hand outstretched. He
stepped back from the advance of these spread fingers.

``Come \gl[hither]{hither}{ici (forme ancienne)},'' said the third blind man, following his
motion and \gl[to clutch]{clutching}{agripper} him neatly.

And they held Nunez and felt him over, saying no word further until they had done so.

``Carefully,'' he cried, with a finger in his eye, and found they thought that organ, with
its fluttering lids, a queer thing in him. They went over it again.

``A strange creature, Correa,'' said the one called Pedro. ``Feel the
\gl[coarseness]{coarseness}{la rudesse, la grossièreté} of his hair. Like a llama's hair.''

``Rough he is as the rocks that \gl[to beget, begot]{begot}{engendrer} him,'' said Correa,
investigating Nunez's unshaven chin with a soft and slightly moist hand. ``Perhaps he will
grow finer.'' Nunez struggled a little under their examination, but they gripped him firm.

``He speaks,'' said the third man. ``Certainly he is a man.''

``And you have come into the world?'' asked Pedro.

``\emph{Out} of the world. Over mountains and glaciers; right over above there, half-way to
the sun. Out of the great big world that goes down, twelve days' journey to the sea.''

They scarcely seemed to \gl[to heed]{heed}{prêter attention à} him. ``Our fathers have told
us men may be made by the forces of Nature,'' said Correa. ``It is the warmth of things and
moisture, and \gl[rottenness]{rottenness}{la pourriture} --- rottenness.''

``Let us lead him to the elders,'' said Pedro.

So they shouted, and Pedro went first and took Nunez by the hand to lead him to the houses.

He drew his hand away. ``I can see,'' he said.

``See?'' said Correa.

``Yes, see,'' said Nunez, turning towards him, and stumbled against Pedro's
\gl[a pail]{pail}{un seau}.

``His senses are still imperfect,'' said the third blind man. ``He stumbles, and talks
unmeaning words. Lead him by the hand.''

``As you will,'' said Nunez, and was led along, laughing.

It seemed they knew nothing of sight.

Well, all in good time he would teach them.
\end{texte}
\source{H.\,G. Wells, \og The Country of the Blind \fg, \emph{The Strand Magazine}, 1904.}

\partiel[12 min]{Questions}
\consigne{\en{Answer in English, in full sentences.}}

\begin{enumerate}[label=\arabic*.,itemsep=2mm,leftmargin=*]
  \item \en{How do the three men judge Nunez before he speaks?}
    \lignes{2}
    \corr{By his unfamiliar steps: they stand with their ears turned towards him instead of
    looking at him. Later they judge him by feeling his hair, his chin and his coat.}
  \item \en{What does the proverb promise Nunez, and what does the story do with it?}
    \lignes{3}
    \corr{It promises that a man who can see will rule those who cannot. The rest of the
    scene contradicts it at once: his sight makes him stumble, his words are taken for
    nonsense, and he is led by the hand like a newborn.}
  \item \en{How do the blind men explain where Nunez comes from?}
    \lignes{3}
    \corr{They decide he came out of the rocks and was made by the forces of Nature ---
    warmth, moisture and rottenness. Everything he says that does not fit is treated as proof
    that his mind is not yet formed.}
  \item \en{Compare their senses with his. Who is better equipped in this valley?}
    \lignes{3}
    \corr{They hear his steps, feel the texture of his hair and skin, and move about a room
    black as pitch without hesitating. He stumbles against a pail. In their world his eyes
    are useless and their hands are exact.}
  \item \en{``It seemed they knew nothing of sight. Well, all in good time he would teach
    them.'' What do these two sentences tell us about Nunez?}
    \lignes{3}
    \corr{That he has understood nothing. He still believes the proverb, and reads his own
    helplessness as their ignorance. The narrator lets him think so, which is how the irony
    works.}
  \item \en{Find two details that show the villagers are not hostile.}
    \lignes{2}
    \corr{Their expression is near awe, not anger; they call it \emph{a marvellous occasion};
    they shout first so that the children will not be frightened; they lead him by the hand
    to help him.}
\end{enumerate}

%% ===================================================================
\newpage
\section{Version}
\consigne{Traduisez le passage en français. Il suit immédiatement le texte.}

\begin{version}
He found it tax his nerve and patience more than he had anticipated, that first encounter with
the population of the Country of the Blind. The place seemed larger as he drew near to it, and
the \gl[smeared]{smeared}{barbouillé} \gl[plastering]{plasterings}{le crépi, l'enduit} queerer,
and a crowd of children and men and women (the women and girls, he was pleased to note, had
some of them quite sweet faces, for all that their eyes were shut and sunken) came about him,
holding on to him, touching him with soft, sensitive hands, smelling at him, and listening at
every word he spoke. Some of the maidens and children, however, kept
\gl[aloof]{aloof}{à l'écart} as if afraid, and indeed his voice seemed coarse and rude beside
their softer notes. They \gl[to mob]{mobbed}{assaillir en foule} him. His three guides kept
close to him with an effect of \gl[proprietorship]{proprietorship}{un air de propriétaire},
and said again and again, ``A wild man out of the rock.''

``Bogota,'' he said. ``Bogota. Over the mountain crests.''

``A wild man --- using wild words,'' said Pedro. ``Did you hear that --- \emph{Bogota}? His
mind is hardly formed yet. He has only the beginnings of speech.''

A little boy \gl[to nip]{nipped}{pincer} his hand. ``Bogota!'' he said mockingly.

``Ay! A city to your village. I come from the great world --- where men have eyes and see.''

``His name's Bogota,'' they said.
\end{version}
\source{\emph{Ibid.}}

\lignes{20}

\corr{\textbf{Proposition de traduction.} Cette première rencontre avec la population du Pays
des Aveugles mit ses nerfs et sa patience à plus rude épreuve qu'il ne s'y était attendu. Le
village paraissait plus grand à mesure qu'il en approchait, et les enduits barbouillés plus
étranges ; et une foule d'enfants, d'hommes et de femmes (les femmes et les filles, nota-t-il
avec plaisir, avaient pour certaines des visages fort doux, malgré leurs yeux clos et
enfoncés) se pressa autour de lui, s'accrochant à lui, le touchant de mains douces et
sensibles, le flairant, écoutant chacun de ses mots. Quelques jeunes filles et quelques
enfants, pourtant, se tenaient à l'écart comme effrayés, et de fait sa voix semblait rude et
grossière à côté de leurs intonations plus douces. Ils l'assaillirent en foule. Ses trois
guides demeuraient près de lui d'un air de propriétaires, et répétaient sans cesse : \og Un
homme sauvage, sorti du rocher. \fg

\og Bogota, dit-il. Bogota. Par-dessus les crêtes. \fg

\og Un homme sauvage --- qui emploie des mots sauvages, dit Pedro. Avez-vous entendu ça ?
\emph{Bogota} ? Son esprit est à peine formé. Il n'a que les commencements de la parole. \fg

Un petit garçon lui pinça la main. \og Bogota ! \fg{} fit-il d'un ton moqueur.

\og Ah ! Une ville, en comparaison de votre village. Je viens du grand monde --- là où les
hommes ont des yeux et voient. \fg

\og Il s'appelle Bogota \fg, dirent-ils.}

\note[Choix de traduction 1]{\emph{He found it tax his nerve and patience\ldots, that first
encounter} : le sujet réel est rejeté à la fin et annoncé par \emph{it}. Le français remonte
le sujet en tête. Noter aussi \emph{tax} sans \emph{to} : construction \emph{find + complément
+ infinitif nu}.}

\note[Choix de traduction 2]{\emph{The place seemed larger as he drew near} : \emph{as} de
progression, à rendre par \og à mesure que \fg, non par \og comme \fg{} ni \og quand \fg.}

\note[Choix de traduction 3]{\emph{for all that their eyes were shut} : \emph{for all that} =
\og malgré \fg, \og en dépit de \fg. Locution figée à ne pas décomposer.}

\note[Choix de traduction 4]{\emph{his voice seemed coarse and rude beside their softer
notes} : \emph{beside} = \og à côté de \fg{} (comparaison), à ne pas confondre avec
\emph{besides} = \og en outre \fg. Un \emph{s} change tout.}

\note[Choix de traduction 5]{\emph{A city to your village} : \emph{to} marque ici la
comparaison, comme dans \emph{nothing to what I saw}. Le français doit étoffer : \og une
ville, en comparaison de votre village \fg.}

%% ===================================================================
\partiel{Lexique à mémoriser}
\begin{itemize}[label=--,itemsep=0.8mm,leftmargin=*]\small
  \item \textbf{blind} — aveugle ; \textit{nom} blindness\textit{, et} a blind \textit{« un store »}
  \item \textbf{sight} — la vue ; out of sight\textit{, at first sight}
  \item \textbf{an eyelid} — une paupière ; eye + lid \textit{« couvercle »}
  \item \textbf{sunken} — enfoncé ; \textit{participe ancien de} to sink
  \item \textbf{awe} — l'effroi mêlé de respect ; \textit{adjectif} awful \textit{« affreux », d'abord « qui inspire l'effroi »}
  \item \textbf{a crest} — une crête
  \item \textbf{coarse} — rude, grossier ; \textit{nom} coarseness
  \item \textbf{rough} — rugueux, brutal ; \textit{se prononce} ruf
  \item \textbf{moist} — humide ; \textit{nom} moisture
  \item \textbf{to shrink} — rétrécir ; shrank, shrunk
  \item \textbf{to startle} — faire sursauter
  \item \textbf{to clutch} — agripper
  \item \textbf{to grip} — empoigner \textit{(fiche 9)}
  \item \textbf{to stumble} — trébucher
  \item \textbf{to heed} — prêter attention à ; \textit{d'où} heedless \textit{« insouciant »}
  \item \textbf{an elder} — un ancien ; \textit{comparatif ancien de} old
  \item \textbf{a maiden} — une jeune fille \textit{(littéraire)}
  \item \textbf{aloof} — à l'écart, distant
  \item \textbf{mockingly} — d'un ton moqueur ; \textit{verbe} to mock
  \item \textbf{all in good time} — chaque chose en son temps
\end{itemize}

\end{document}
